\section{Universality beyond Gaussian smoothing}
Gaussian smoothing makes several formulas explicit, but the logarithmic
threshold and the sharp fixed-codimension constant do not require Gaussian
noise. Let $V$ have a symmetric absolutely continuous law with
$\E V^2=1$ and $\E V^4<\infty$. The noise law is fixed as dimension
grows. Set $W=\xi+\sigma V$, with independent Rademacher $\xi$, and
write $\varphi(u)=\E\cos(uV)$. In particular this includes standardized
uniform and Laplace smoothing. Uniformly perturbed Rademacher
initializations also appear in the experiments of Abbe et al.; their use
as an initialization scheme is not a new proposal here.

\begin{theorem}[A noise-law universal logarithmic transition]
For constants $c,C,\sigma_0>0$ depending only on the law of $V$, the
zero-bias full-parity feature eigenvalue satisfies
\[
 c r^{-1}e^{-C\sigma^2r}\le\lambda_r
       \le C e^{-c\sigma^2r},\qquad 0\le\sigma\le\sigma_0.
 \tag{30}
\]
The upper bound also holds for $\sigma_0\le\sigma\le1$, after changing
the constants. Thus the polynomial/superpolynomial boundary remains
$\sigma^2r$ of order $\log r$.

More precisely, fix $A>0$ and a nonnegative integer $k$, let $r$ tend
to infinity through even integers, and put $\sigma^2=A\log r/r$.
For any target omitting $2k$ coordinates, its eigenvalue obeys
\[
 \lambda_{r,r-2k}\sim
 \frac{32}{\pi^5}\left(\frac\pi2\right)^{4k}A^{2k}
 (\log r)^{2k}r^{-2k-1-A\pi^2/4}.                    \tag{31}
\]
The leading constant is the same as for Gaussian smoothing of variance
one. The noise law must remain fixed; no uniformity over changing laws
or their fourth moments is asserted.
\end{theorem}
\begin{proof}
Absolute continuity and the expansion at zero imply
\[
 |\varphi(u)|\le1-c_1\min\{u^2,1\}\quad(u\in\R)
 \tag{32}
\]
for some $c_1>0$. Indeed $\varphi(u)=1-u^2/2+O(u^4)$ near zero;
on a compact set away from zero its modulus is strictly below one,
and it tends to zero at infinity. A density cannot be supported on the
arithmetic set required for equality $|\varphi(u)|=1$ at $u\ne0$.
For $|t|\le1/4$, $\E\sin^2(tW)\le(1+\sigma^2)t^2\le1/8$.
For $|t|\ge1/4$, (32) gives
\[
 s_\sigma(t):=\E\sin^2(tW)
       =\frac{1-\cos(2t)\varphi(2\sigma t)}2
       \le1-c_2\sigma^2\qquad(0<\sigma\le1).
\]
Equation (2) proves the upper bound in (30). For sufficiently small
$\sigma_0$, $\varphi(\sigma t)\ge\exp(-C_1\sigma^2t^2)$ when
$0\le t\le\pi/2$. The lower witness (5) now has integrand
$\varphi(\sigma t)^r\sin^r(t)/t^2$, which is nonnegative because
$r$ is even. The same first-peak interval as before proves the lower
bound in (30). This part only needs a finite second moment.

For the sharp assertion let $j=r-2k$ and define
\[
 J_{r,k}=\frac1\pi\int_0^\infty
       s_\sigma(t)^{j/2}(1-s_\sigma(t))^k\frac{dt}{t^2}.
\]
Minkowski gives $\sqrt{\lambda_{r,j}}\le J_{r,k}$ as in the Gaussian
case. For a lower witness, gauge the latent signs so that the incoming
weights become $1+\sigma V_i$, and multiply the Fourier coefficient by
$Q=\prod_{i\notin T}V_i$. Independence gives $\E Q^2=1$, while
$\E[V\cos(t+\sigma tV)]=\sin t\,\varphi'(\sigma t)$.
The cosine-integral representation therefore gives
\[
 \sqrt{\lambda_{r,j}}\ge\nu_{r,k}:=
 \frac1\pi\int_0^\infty
 \frac{\sin^r t\,\varphi(\sigma t)^j
                    \varphi'(\sigma t)^{2k}}{t^2}\,dt.
 \tag{33}
\]
The displayed integrand is nonnegative. The witness is a product of
first powers of the noise coordinates, not an assertion that the noise
has Gaussian Hermite polynomials. Fubini for truncated integrals and
Cauchy--Schwarz with $\E Q^2\|w\|^2<\infty$ justify the identity;
the fourth moment is enough for this latter bound.

Put $t_0=\pi/2$ and $L=\sigma^2r=A\log r$.
In a shrinking neighborhood $t=t_0+u/\sqrt r$ with bounded $u$,
\[
 \varphi(\sigma t)^j\sin^r t
     =e^{-Lt_0^2/2-u^2/2}(1+o(1)),\qquad
 \varphi'(\sigma t)^{2k}=(\sigma t_0)^{2k}(1+o(1)).
\]
For the upper integrand one has, on the same window,
\[
 s_\sigma(t)^{j/2}
     =e^{-Lt_0^2/2-u^2/2}(1+o(1)),\qquad
 (1-s_\sigma(t))^k=(\sigma^2t_0^2)^k(1+o(1)).
\]
The estimates use $r\sigma^4=O((\log r)^2/r)=o(1)$ and
$L/\sqrt r=o(1)$. The latter factor for $k>0$ is dominated by
$\sigma^2t_0^2$, because $L\to\infty$; this explains the power
$(\log r)^{2k}$ in the squared result.

Here are bounds justifying localization of the upper integral. For
$|\sigma t|\le u_0$, with $u_0$ small and fixed, the identity
\[
 1-s_\sigma(t)=\varphi(2\sigma t)\cos^2t+
                         \{1-\varphi(2\sigma t)\}/2
\]
gives $1-s_\sigma(t)\ge c_3(\cos^2t+\sigma^2t^2)$.
On each interval around $t_m=(2m+1)\pi/2$ it follows that the upper
integrand has a Gaussian envelope in $\sqrt r(t-t_m)$ and a factor
$e^{-c_4Lt_m^2}$. The additional factor $(1-s_\sigma)^k$ is bounded
there by a constant times $\{(t-t_m)^2+\sigma^2t_m^2\}^k$.
On a fixed small interval near zero use instead
$s_\sigma(t)\le2t^2$; its integral is exponentially small in $r$.
When $|\sigma t|>u_0$, (32) gives a fixed gap from one, so the
remaining $t^{-2}$ integral is $O(\sigma e^{-c_5r})$.
For the intervening peaks choose a fixed large $M$ so that
$c_4t_M^2>t_0^2/2+1/A$. Their summed envelope is negligible relative
to the first peak. Each of the finitely many peaks $1\le m<M$ has
the same local expansion with $t_0$ replaced by $t_m$, hence is smaller
by $r^{-A(t_m^2-t_0^2)/2}$ up to a constant depending on $m,k$.
These estimates also dominate expanding first-peak windows; their
Gaussian tails permit the bounded $u$ cutoff to tend to infinity.

Thus the upper integral and the first-peak part of the nonnegative lower
integral have the same asymptotic value
\[
 J_{r,k}\sim
 \frac{\sqrt{2\pi}}\pi\,t_0^{2k-2}\sigma^{2k}r^{-1/2}
                                      e^{-Lt_0^2/2}.
\]
Squeezing between (33) and $J_{r,k}$ and squaring proves (31).
\end{proof}

\paragraph{Finite features without a sub-Gaussian noise assumption.}
The positive full-parity statement also supplies a polynomial-width
feature map under these noise laws. Let $\mu$ be the signed lower witness
in (30), and use the latent-sign comparator (7). Symmetrization and
Rademacher contraction, now only in expectation, give
\[
 \E\sup_{x\in X}|q_N(x)-\chi(x)|
       \le\frac{4r\sqrt{1+\sigma^2}}{|\mu|\sqrt N}.
\]
Indeed conditional Rademacher signs absorb the latent Boolean coefficient,
and the supremum of a linear form on the cube is at most $\sqrt r$
times its Euclidean norm. The second moment of each incoming vector is
$r(1+\sigma^2)$. Markov's inequality shows that
$N\ge256r^2(1+\sigma^2)/(\mu^2\zeta^2)$ suffices for uniform
comparator error at most $1/2$, except with probability $\zeta/2$.
Also the empirical feature radius is at most
$2r^2(1+\sigma^2)/\zeta$ except with probability $\zeta/2$.
The comparator norm is exactly at most $|\mu|^{-1}$, so a perceptron
has at most $8r^2(1+\sigma^2)/(\zeta\mu^2)$ mistakes on every input
sequence on the joint event. Confidence dependence is polynomial here;
no sub-Gaussian tail was assumed.

For fixed codimension, use instead the coefficient
$\chi_T(\xi)Q(V)$ and the nonzero signed witness (33). The same argument
uses $\E Q^2=1$ and
$\E[Q^2\|w\|^2]\le C_{k,V}r(1+\sigma^2)$, while an additional
Markov bound controls the empirical coefficient norm. Equation (31)
therefore gives polynomial sufficient width and arbitrary-sequence
perceptron mistakes for every fixed $k$ at $\sigma^2r=A\log r$.
These are frozen-feature guarantees; no additional full-network
optimization theorem is inferred from the noise-law extension.
